% ======================================================================
% 3.4 Оптимизация затрат. Ответственный: A — Плешнев Арсений.
%
% Что сдать (main.tex: теория 1380–1414, порядок 1415–1436,
% разобранный пример 1438–1493; требования к отчёту 1871–1879):
%   1. Объяснить, почему в минимуме производная по любому направлению равна
%      нулю (grad Phi = 0); найти стационарную точку, выразить c_i* через
%      базовую нагрузку R_i=lambda/mu_i (правило квадратного корня).
%   2. Доказать, что точка — минимум (Гессе, выпуклость -> глобальный минимум);
%      вычислить c1* и сравнить затраты первого звена при исходном c и при c1*.
%   3. Найти оптимальную целевую загрузку u0*=lambda/(mu1*c1*) и соответствующий
%      порог k* из п. 3.3.
%   4. При ограничении c1+c2=N: условие оптимальности через производную
%      по направлению (1,-1) и то же условие методом Лагранжа; найти c1,c2
%      численно; проверить сумму и положительность запасов; смысл множителя.
%
% Параметры варианта 3 (приложение Б): lambda=600, mu1=mu=80, mu2=120,
% a=1, b=1, N=16; исходное c=8 (рабочая точка из п. 3.1).
%
% Программные значения — task_3_4 (python experiments.py 3.4);
% градиентный спуск для этой же функции — п. 3.5.5 (файл 353-simulator.tex, зона C).
% ======================================================================

\subsection{Оптимизация затрат}

\taskhead{Исходные данные и функция затрат}
\todo{A: записать $\Phi(c_1,c_2)$ с параметрами варианта 3; допустимая область $s_i=\mu_ic_i-\lambda>0$; базовые нагрузки $R_i=\lambda/\mu_i$.}

\taskhead{Стационарная точка и доказательство минимума}
\todo{A: объяснить, почему $\operatorname{grad}\Phi=0$ в точке минимума; найти стационарную точку и выразить $c_i^*$ через $R_i$; доказать минимум (вторые производные, выпуклость на всей допустимой области -> глобальный характер); вычислить значения (пример: main.tex, строки 1440–1457).}

\taskhead{Затраты и оптимальные настройки регулятора}
\todo{A: сравнить затраты первого звена при исходном $c$ и при $c_1^*$; найти $u_0^*=\lambda/(\mu_1c_1^*)$ и соответствующий порог $k_*$ по формуле п. 3.3.}

\taskhead{Ограниченный общий ресурс $c_1+c_2=N$}
\todo{A: получить условие оптимального распределения из равенства нулю производной по направлению $(1,-1)$; получить то же условие методом Лагранжа; найти $c_1,c_2$ численно (python experiments.py 3.4); проверить ограничение на сумму и положительность запасов; пояснить смысл множителя Лагранжа (пример: main.tex, строки 1467–1493).}

\taskhead{Сопоставление с программным расчётом}
% Программные значения — вызовы своих функций из solution.py (task_3_4).
\todo{A: листинг вызовов функций и запуск python experiments.py 3.4.}
\begin{table}[htbp]
\centering\footnotesize
\setlength{\tabcolsep}{4pt}
\caption{Сравнение аналитических и программных значений (п. 3.4)}
\label{tab:3-4}
\begin{tabular}{|l|r|r|r|r|}\hline
Величина & \shortstack{Аналитическое\\значение} & \shortstack{Программное\\значение} & \shortstack{Абсолютная\\ошибка} & \shortstack{Относительная\\ошибка}\\\hline
$c_1^*$ & — & — & — & —\\
$c_2^*$ & — & — & — & —\\
$\Phi^*$ & — & — & — & —\\
$u_0^*$ & — & — & — & —\\
$k_*$ при $u_0^*$ & — & — & — & —\\
$c_1$ при $c_1+c_2=N$ & — & — & — & —\\
$c_2$ при $c_1+c_2=N$ & — & — & — & —\\
\hline\end{tabular}
\end{table}
% Для распределения при ограничении N аналитического ответа в замкнутой форме
% может не быть: тогда указать способ проверки (main.tex, строки 1911–1915)
% и не называть разность двух приближений ошибкой относительно точного решения.

\taskhead{Вывод}
\todo{A: правило квадратного корня и его смысл; выигрыш в затратах; эффект ограничения ресурса.}
